% GENERATED FILE. Edit canonical lessons, then run npm run generate:books.
% canonical-source-hash: fnv1a64:7eef9e5f54b02563
% canonical-lessons: UR-C83-anguthi, UR-C83-ghari, UR-C83-gyarah, UR-C83-barah, UR-C83-terah

\chapter{Ring, Watch, Eleven, Twelve, Thirteen}
\label{ch:ur-ring-watch-eleven-twelve-thirteen}
\hlchaptermodality{83}

\begin{chapteropening}
\textbf{By the end of this chapter:} \emph{I can say a ring, a watch, a clock, eleven, twelve and thirteen.}

Complete the last lesson of chapter 83: I can say a ring, a watch, a clock, eleven, twelve and thirteen.
\end{chapteropening}

\section[\texorpdfstring{\ur{انگوٹھی} (angūṭhī)}{انگوٹھی (angūṭhī)}]{\ur{انگوٹھی} (angūṭhī) — a ring}

\label{lesson:UR-C83-anguthi}



\begin{quote}
Before the new one: say the Urdu for a sock, then the Urdu for a pocket.
\end{quote}

\subsection*{You'll want to know: \ur{انگوٹھی}}
\textbf{\ur{انگوٹھی}} — \emph{angūṭhī} — \textquotedblleft{}a ring\textquotedblright{}.

\begin{cousinweb}
\textbf{\ur{انگوٹھی}} (\emph{angūṭhī}) is a tadbhava, a word inherited and reshaped by sound change. It comes from \emph{*aṅguṣṭhiya}, \textquotedblleft{}of the thumb\textquotedblright{}, a form CDIAL reconstructs from its descendants (CDIAL 138). It is built on \emph{aṅguṣṭha}, \textquotedblleft{}thumb\textquotedblright{}.
\end{cousinweb}

\begin{grammarlens}[title={What you've built}]
One more word for this chapter.
\end{grammarlens}

\subsection*{Your turn}
\begin{itemize}
\raggedright
  \item \emph{Say it:} \emph{angūṭhī}
  \item \emph{Say it:} \emph{angūṭhī}, once more
  \item \emph{Say it:} say \emph{angūṭhī}
  \item \emph{Recall:} say the Urdu for a sock, then the Urdu for a pocket, then say \emph{angūṭhī} again
\end{itemize}

\begin{tcolorbox}[breakable,colback=teal!4,colframe=teal!35!black,title={Before you move on}]
What does \ur{انگوٹھی} mean? (\textquotedblleft{}A ring\textquotedblright{}.) Say it once more.
\end{tcolorbox}
\section[\texorpdfstring{\ur{گھڑی} (ghaṛī)}{گھڑی (ghaṛī)}]{\ur{گھڑی} (ghaṛī) — a watch, a clock}

\label{lesson:UR-C83-ghari}



\begin{quote}
Before the new one: say the Urdu for a pocket, then the Urdu for a ring.
\end{quote}

\subsection*{You'll want to know: \ur{گھڑی}}
\textbf{\ur{گھڑی}} — \emph{ghaṛī} — \textquotedblleft{}a watch, a clock\textquotedblright{}.

\begin{cousinweb}
\textbf{\ur{گھڑی}} (\emph{ghaṛī}) is a tadbhava, a word inherited from Sanskrit and reshaped by sound change. It comes from \emph{ghaṭī}, \textquotedblleft{}pot, water-clock\textquotedblright{} (CDIAL 4406). A \emph{ghaṭī} was a small pot that sank as it filled, timing a short period.
\end{cousinweb}

\begin{grammarlens}[title={What you've built}]
One more word for this chapter.
\end{grammarlens}

\subsection*{Your turn}
\begin{itemize}
\raggedright
  \item \emph{Say it:} \emph{ghaṛī}
  \item \emph{Say it:} \emph{ghaṛī}, once more
  \item \emph{Say it:} say \emph{ghaṛī}
  \item \emph{Recall:} say the Urdu for a pocket, then the Urdu for a ring, then say \emph{ghaṛī} again
\end{itemize}

\begin{tcolorbox}[breakable,colback=teal!4,colframe=teal!35!black,title={Before you move on}]
What does \ur{گھڑی} mean? (\textquotedblleft{}A watch, a clock\textquotedblright{}.) Say it once more.
\end{tcolorbox}
\section[\texorpdfstring{\ur{گیارہ} (gyārah)}{گیارہ (gyārah)}]{\ur{گیارہ} (gyārah) — eleven}

\label{lesson:UR-C83-gyarah}



\begin{quote}
Before the new one: say the Urdu for a ring, then the Urdu for a watch.
\end{quote}

\subsection*{You'll want to know: \ur{گیارہ}}
\textbf{\ur{گیارہ}} — \emph{gyārah} — \textquotedblleft{}eleven\textquotedblright{}.

\begin{cousinweb}
\textbf{\ur{گیارہ}} (\emph{gyārah}) is a tadbhava, a word inherited from Sanskrit and reshaped by sound change. It comes from \emph{ekādaśa}, \textquotedblleft{}eleven\textquotedblright{} (CDIAL 2485).
\end{cousinweb}

\begin{grammarlens}[title={What you've built}]
One more word for this chapter.
\end{grammarlens}

\subsection*{Your turn}
\begin{itemize}
\raggedright
  \item \emph{Say it:} \emph{gyārah}
  \item \emph{Say it:} \emph{gyārah}, once more
  \item \emph{Say it:} say \emph{gyārah}
  \item \emph{Recall:} say the Urdu for a ring, then the Urdu for a watch, then say \emph{gyārah} again
\end{itemize}

\begin{tcolorbox}[breakable,colback=teal!4,colframe=teal!35!black,title={Before you move on}]
What does \ur{گیارہ} mean? (\textquotedblleft{}Eleven\textquotedblright{}.) Say it once more.
\end{tcolorbox}
\section[\texorpdfstring{\ur{بارہ} (bārah)}{بارہ (bārah)}]{\ur{بارہ} (bārah) — twelve}

\label{lesson:UR-C83-barah}



\begin{quote}
Before the new one: say the Urdu for a watch, then the Urdu for eleven.
\end{quote}

\subsection*{You'll want to know: \ur{بارہ}}
\textbf{\ur{بارہ}} — \emph{bārah} — \textquotedblleft{}twelve\textquotedblright{}.

\begin{cousinweb}
\textbf{\ur{بارہ}} (\emph{bārah}) is a tadbhava, a word inherited from Sanskrit and reshaped by sound change. It comes from \emph{dvādaśa}, \textquotedblleft{}twelve\textquotedblright{} (CDIAL 6658).
\end{cousinweb}

\begin{grammarlens}[title={What you've built}]
One more word for this chapter.
\end{grammarlens}

\subsection*{Your turn}
\begin{itemize}
\raggedright
  \item \emph{Say it:} \emph{bārah}
  \item \emph{Say it:} \emph{bārah}, once more
  \item \emph{Say it:} say \emph{bārah}
  \item \emph{Recall:} say the Urdu for a watch, then the Urdu for eleven, then say \emph{bārah} again
\end{itemize}

\begin{tcolorbox}[breakable,colback=teal!4,colframe=teal!35!black,title={Before you move on}]
What does \ur{بارہ} mean? (\textquotedblleft{}Twelve\textquotedblright{}.) Say it once more.
\end{tcolorbox}
\section[\texorpdfstring{\ur{تیرہ} (terah)}{تیرہ (terah)}]{\ur{تیرہ} (terah) — thirteen}

\label{lesson:UR-C83-terah}



\begin{quote}
Before the new one: say the Urdu for eleven, then the Urdu for twelve.
\end{quote}

\subsection*{You'll want to know: \ur{تیرہ}}
\textbf{\ur{تیرہ}} — \emph{terah} — \textquotedblleft{}thirteen\textquotedblright{}.

\begin{cousinweb}
\textbf{\ur{تیرہ}} (\emph{terah}) is a tadbhava, a word inherited and reshaped by sound change. It comes from \emph{*trayedaśa}, \textquotedblleft{}thirteen\textquotedblright{}, a form CDIAL reconstructs from its descendants (CDIAL 6001). It is a variant CDIAL reconstructs beside Sanskrit \emph{trayodaśa}, \textquotedblleft{}thirteen\textquotedblright{}.
\end{cousinweb}

\begin{grammarlens}[title={What you've built}]
One more word for this chapter.
\end{grammarlens}

\subsection*{Your turn}
\begin{itemize}
\raggedright
  \item \emph{Say it:} \emph{terah}
  \item \emph{Say it:} \emph{terah}, once more
  \item \emph{Say it:} say \emph{terah}
  \item \emph{Recall:} say the Urdu for eleven, then the Urdu for twelve, then say \emph{terah} again
\end{itemize}

\begin{tcolorbox}[breakable,colback=teal!4,colframe=teal!35!black,title={Before you move on}]
What does \ur{تیرہ} mean? (\textquotedblleft{}Thirteen\textquotedblright{}.) Say it once more.
\end{tcolorbox}
